\section{A collision-regular integer hierarchy}
\label{sec:integer}

On the tube \eqref{eq:domain}, set \(w_j=\rho_j-c\), \(Q_j=w_j^2-D\),
and \(m_j=Q_j^{-1}\) for each exterior occurrence. The rational
moments are
\[
 S_n=\sum_j m_j^n\quad(n\geq 1),\qquad
 T_n=\sum_j w_jm_j^n\quad(n\geq 2),\qquad M=S_1 .
\]
They are real by conjugation. The collar gives
\(\lvert Q_j\rvert\geq(15/16)\lvert w_j\rvert^2\). This bound and the
unit count in \eqref{eq:count-safe} give absolute convergence in the
indicated index ranges. In particular, the first odd moment is outside
these ranges. The hierarchy extends the interaction law
of \cite{BakerGT41} and the next even moment calculation of
\cite{BakerGR02}, using the rational divided differences of
\cite{BakerCM09}. Primes denote physical time derivatives.

\begin{theorem}[Integer moment evolution]\label{thm:hierarchy}
On the complete canonical tube, the regular moments satisfy
\begin{align}
 S_n'={}&2n\sum_{k=1}^{n}S_kS_{n+1-k}
       +2n\sum_{k=2}^{n}\bigl(T_kT_{n+2-k}+DS_kS_{n+2-k}\bigr)
       \notag\\[-1pt]
       &\quad-4n(n+2)S_{n+1}-4n(n+3)DS_{n+2},
       \qquad n\geq1, \label{eq:S-hierarchy}\\
 T_n'={}&2(2n-1)\sum_{k=2}^{n}T_kS_{n+1-k}
       +4nD\sum_{k=2}^{n}T_kS_{n+2-k}
       \notag\\[-1pt]
       &\quad-4(n^2+n-1)T_{n+1}-4n(n+3)DT_{n+2},
       \qquad n\geq2. \label{eq:T-hierarchy}
\end{align}
Empty sums are zero. These equations and \(D'=2-4DM\) hold through
target and exterior multiplicities, with inward derivatives at the
interval endpoints. Each equation involves finitely many terms of the
infinite relative moment family. The absolute midpoint \(c\) uses
its counted contour law.
\end{theorem}

\begin{proof}
First consider a polynomial satisfying the physical heat equation at
a slice with simple roots. Factoring out a simple root gives
\(\rho'=2\sum_{\eta\ne\rho}(\rho-\eta)^{-1}\). At this finite stage
write \(P=\sum_j w_jm_j=T_1\). Adding the target velocities gives
\(c'=-2P\). Multiplying their velocity difference by \(s\), where
\(s^2=D\), gives \(D'\). The mutual target force contributes \(2\)
to this derivative, and each exterior occurrence contributes
\(-4Dm_j\). Thus \(D'=2-4DM\), as in \eqref{eq:pairlaw}.
The contribution of the targets to the velocity at an exterior root
is \(4w_jm_j\). Hence
\begin{equation}
 w_j'=4w_jm_j+2P+2\sum_{k\ne j}\frac{1}{w_j-w_k},
 \label{int:relative}
\end{equation}
where the last sum ranges over exterior occurrences.

For the symmetrization, let \(A=(u^2-D)^{-1}\), \(B=(v^2-D)^{-1}\),
and \(L_h=\sum_{i=0}^{h-1}A^{h-1-i}B^i\) for \(h\geq1\), with
\(L_0=0\). The two divided differences are
\begin{align}
 \mathcal E_h(u,v)&=\frac{A^h-B^h}{u-v}=-(u+v)AB L_h, \notag\\
 \mathcal O_h(u,v)&=\frac{uA^h-vB^h}{u-v}
              =-AB L_{h-1}-(D+uv)AB L_h. \label{int:divdiff}
\end{align}
Indeed \(A-B=-(u-v)(u+v)AB\), and
\(uA-vB=-(u-v)(D+uv)AB\). The identity
\(x^h=xL_h-AB L_{h-1}\) for \(x=A,B\) follows by cancellation of
successive terms. It gives the second line of \eqref{int:divdiff}
from the first-power difference. Both rational expressions on the
right are regular at \(u=v\) away from zeros of the denominators.

The product \(AB L_h\) is \(\sum_{k=1}^h A^k B^{h+1-k}\).
The symmetric kernels in \eqref{int:divdiff} therefore give
\begin{align}
 -\sum_{i<j}\mathcal O_{n+1}(w_i,w_j)
  ={}&\frac12\sum_{k=1}^{n}
             \bigl(S_kS_{n+1-k}-S_{n+1}\bigr)\notag\\
   &+\frac D2\sum_{k=1}^{n+1}
             \bigl(S_kS_{n+2-k}-S_{n+2}\bigr)\notag\\
   &+\frac12\sum_{k=1}^{n+1}
             \bigl(T_kT_{n+2-k}-S_{n+1}-DS_{n+2}\bigr),
     \label{int:odd-pairs}\\
 -\sum_{i<j}\mathcal E_h(w_i,w_j)
  ={}&\sum_{k=1}^{h}\bigl(T_kS_{h+1-k}-T_{h+1}\bigr).
     \label{int:even-pairs}
\end{align}
The factor \(1/2\) in \eqref{int:odd-pairs} converts the ordered
off-diagonal sum to the unordered sum. The three diagonal subtractions
come from \(A^{n+1}\), \(DA^{n+2}\), and
\(w^2 A^{n+2}=A^{n+1}+DA^{n+2}\), respectively.
The two contributions from \(u+v\) in the even difference give the
same ordered sum after \(k\) is reversed. Each diagonal contributes
\(T_{h+1}\), and the two terms cancel the ordering factor \(1/2\).

For the even weight \(f(w,D)=(w^2-D)^{-n}\),
\(f_D=n(w^2-D)^{-n-1}\) and \(f_w=-2nw(w^2-D)^{-n-1}\).
Differentiating along \eqref{int:relative} gives
\begin{align}
 S_n'={}&nD'S_{n+1}-8n(S_{n+1}+DS_{n+2})
          -4nPT_{n+1}-4n\sum_{i<j}\mathcal O_{n+1}(w_i,w_j).
 \label{int:even-before}
\end{align}
The \(k=1,n+1\) terms of the \(T\)-convolution in
\eqref{int:odd-pairs} contribute \(4nPT_{n+1}\) in
\eqref{int:even-before}. The same two endpoints of the \(D\)-weighted
convolution contribute \(4nDMS_{n+1}\), which cancels the interaction
term of \(nD'S_{n+1}\). The remaining coefficient of \(S_{n+1}\)
is \(2n-8n-2n^2-2n(n+1)=-4n(n+2)\). The coefficient of
\(DS_{n+2}\) is \(-8n-4n(n+1)=-4n(n+3)\). This yields
\eqref{eq:S-hierarchy}.

For the odd weight \(f(w,D)=w(w^2-D)^{-n}\) one has
\[
 f_D=nw(w^2-D)^{-n-1},\qquad
 f_w=(1-2n)(w^2-D)^{-n}-2nD(w^2-D)^{-n-1}.
\]
Consequently
\begin{align}
 T_n'={}&nD'T_{n+1}+4(1-2n)T_{n+1}-8nDT_{n+2}\notag\\
       &+2P\bigl((1-2n)S_n-2nDS_{n+1}\bigr)\notag\\
       &+2(1-2n)\sum_{i<j}\mathcal E_n(w_i,w_j)
          -4nD\sum_{i<j}\mathcal E_{n+1}(w_i,w_j).
       \label{int:odd-before}
\end{align}
By \eqref{int:even-pairs}, the \(k=1\) terms cancel both terms
proportional to \(P\). The \(k=n+1\) term of the \(D\)-weighted
sum yields \(4nDMT_{n+1}\), canceling the interaction part of
\(nD'T_{n+1}\). The other diagonal and direct terms give
\(2n+4(1-2n)+2n(1-2n)=-4(n^2+n-1)\) for \(T_{n+1}\) and
\(-8n-4n(n+1)=-4n(n+3)\) for \(DT_{n+2}\).
The two surviving convolution ranges are \(2\leq k\leq n\).
This proves \eqref{eq:T-hierarchy} at the finite simple slice.
The arguments have eliminated every use of \(P\).

For the complete passage, fix a reference slice \(\tau\) and a
stationary positive physical circle \(\Gamma\) enclosing the two
targets and excluding every exterior occurrence. Use
\(u=z-c(t)\), \(Q_*(t,z)=u^2-D(t)\).
For a weight \(f=Q_*^{-n}\) or \(f=uQ_*^{-n}\),
Proposition~\ref{prop:trace} writes the exterior sum as
\(-\langle f(H_z/H-Q_{*,z}/Q_*)\rangle_\Gamma\), where
\(\langle g\rangle_\Gamma=(2\pi\ii)^{-1}\oint_\Gamma g(z)\,\dd z\).
Both poles remain inside the fixed circle, including when they
coalesce. The combined target term integrates to zero since
\(fQ_{*,z}/Q_*=O(z^{-3})\) at infinity and all its poles are inside.
It vanishes as an identity in time. Differentiation of the sum
retains the full motion of the kernel, using \(u_t=-c'\) and
\(Q_{*,t}=-2uc'-D'\), together with
\[
 \partial_t(H_z/H)=-H_{zzz}/H+H_zH_{zz}/H^2 .
\]
The contour expression is locally holomorphic in time.

For a finite multiple-root slice, perturb the root factors at that
slice and initialize a separate polynomial heat solution for each
perturbation. The strict clearance keeps the target count. The
spatial derivatives through order three converge on \(\Gamma\).
The displayed heat derivative and the counted target moments then
give convergence of the first contour jets, including \(c',D'\).
The right sides of \eqref{eq:S-hierarchy} and \eqref{eq:T-hierarchy}
are continuous finite sums of occurrence values, so the identities
persist at polynomial multiples.

For the canonical field, apply Lemma~\ref{lem:slice-passage} at
\(\tau\) with modulus cutoff \(N\), keeping the entire opposite and
conjugate groups and every tied occurrence. The truncated slice
contains the exact targets for large \(N\). Its heat derivative at
\(\tau\) is determined by its first three spatial derivatives on
\(\Gamma\). The omitted inverse-square product tail
\(O(\log N/N)\) gives uniform compact convergence of those
derivatives by Cauchy's formula. Together with the strict clearance
on the fixed circles, this gives convergence of the time jets of
the rational moments and of \(c,D\) at the reference slice.

For \(N>4T+8\), an omitted root has \(\lvert w_j\rvert\geq
\lvert\rho_j\rvert/2\). A dyadic shell of radii comparable to
\(2^kN\) contains \(O(\Cu\,2^kN\log(2^kN))\) roots by the strip
and unit count. Summation against the decay powers \(2n\) and
\(2n-1\) gives the respective moment errors
\(O(N^{1-2n}\log N)\) and \(O(N^{2-2n}\log N)\).
Thus all products in each fixed equation converge. A finite
cover by cleared circles gives the jet passage uniformly on the
compact established tube at fixed \(T\) and fixed orders. The argument
proves the canonical equations using these reference-slice jets alone.
\end{proof}

In particular the first forcing moments and the first odd derivative are
\begin{equation}
 \begin{aligned}
 S_1'&=2M^2-12S_2-16DS_3,\\
 S_2'&=8MS_2-32S_3+4DS_2^2-40DS_4+4T_2^2,\\
 T_2'&=(6M+8DS_2)T_2-20T_3-40DT_4 .
 \end{aligned}
 \label{eq:low-orders}
\end{equation}
The spatial term \(4T_2^2\) in \(S_2'\) has its inherited
positive sign \cite{BakerGR02,BakerGR05}. The equations describe
componentwise evolution for the given field. Arbitrary-sequence
well-posedness, finite truncation, and a signed forcing estimate
are separate questions.
